% 基础数学试卷 - 完整示例
% 包含：选择题、填空题、解答题、草稿纸
% 编译：xelatex 02-math-basic.tex

\documentclass{exam-zh}

% ===== 显示答案（教师版）=====
% 取消下面几行的注释可生成教师版
% \examsetup{
%   question/show-answer = true,
%   solution/show-solution = show-stay,
%   question/show-points = true
% }

\begin{document}

% ===== 试卷抬头 =====
\title{2024年秋季期末考试}
\subject{数学}
\maketitle

\information{
  姓名\underline{\hspace{6em}},
  班级\underline{\hspace{6em}},
  学号\underline{\hspace{6em}}
}

\begin{notice}
  \item 本试卷共 3 大题，满分 100 分，考试时间 120 分钟。
  \item 请用黑色签字笔在答题卡上作答，在试卷上作答无效。
  \item 考试结束后，将试卷和答题卡一并交回。
  \item 可以使用计算器。
\end{notice}

\vspace{1em}

% ===== 一、选择题 =====
\section{选择题（本题共 8 小题，每小题 5 分，共 40 分）}

\begin{question}
  已知集合 $A = \{x \mid x^2 - 3x + 2 = 0\}$，$B = \{1, 2, 3\}$，则 $A \cup B = $ \paren[C]
  \begin{choices}
    \item $\{1\}$
    \item $\{1, 2\}$
    \item $\{1, 2, 3\}$
    \item $\{2, 3\}$
  \end{choices}
\end{question}

\begin{question}
  函数 $f(x) = \ln(x + 1)$ 的定义域是\paren[B]
  \begin{choices}
    \item $\mathbb{R}$
    \item $(-1, +\infty)$
    \item $[0, +\infty)$
    \item $(-\infty, -1)$
  \end{choices}
\end{question}

\begin{question}
  若复数 $z = 1 + \iu$（$\iu$ 为虚数单位），则 $|z| = $ \paren[D]
  \begin{choices}
    \item $1$
    \item $\iu$
    \item $2$
    \item $\sqrt{2}$
  \end{choices}
\end{question}

\begin{question}
  已知向量 $\vec{a} = (2, 1)$，$\vec{b} = (-1, 3)$，则 $\vec{a} \cdot \vec{b} = $ \paren[A]
  \begin{choices}
    \item $1$
    \item $-1$
    \item $5$
    \item $-5$
  \end{choices}
\end{question}

\begin{question}
  等比数列 $\{a_n\}$ 中，$a_1 = 2$，$a_3 = 8$，则 $a_5 = $ \paren[C]
  \begin{choices}
    \item $16$
    \item $24$
    \item $32$
    \item $64$
  \end{choices}
\end{question}

\begin{question}
  已知 $\cos \alpha = \frac{3}{5}$，且 $\alpha \in (0, \frac{\uppi}{2})$，则 $\sin \alpha = $ \paren[B]
  \begin{choices}
    \item $\frac{3}{5}$
    \item $\frac{4}{5}$
    \item $\frac{5}{4}$
    \item $\frac{5}{3}$
  \end{choices}
\end{question}

\begin{question}
  圆 $(x - 1)^2 + (y + 2)^2 = 9$ 的圆心坐标和半径分别是\paren[A]
  \begin{choices}
    \item $(1, -2)$，$3$
    \item $(-1, 2)$，$3$
    \item $(1, -2)$，$9$
    \item $(-1, 2)$，$9$
  \end{choices}
\end{question}

\begin{question}
  函数 $f(x) = \eu^x - x$ 的零点所在的区间是\paren[D]
  \begin{choices}
    \item $(-1, 0)$
    \item $(1, 2)$
    \item $(2, 3)$
    \item $(0, 1)$
  \end{choices}
\end{question}


% ===== 二、填空题 =====
\section{填空题（本题共 4 小题，每小题 5 分，共 20 分）}
\examsetup{question/index = 0}

\begin{question}
  若 $f(x) = 2x^2 - 3x + 1$，则 $f(1) = $ \fillin[0]。
\end{question}

\begin{question}
  已知 $\tan \alpha = 2$，则 $\dfrac{\sin \alpha + \cos \alpha}{\sin \alpha - \cos \alpha} = $ \fillin[3]。
\end{question}

\begin{question}
  双曲线 $\dfrac{x^2}{9} - \dfrac{y^2}{16} = 1$ 的渐近线方程为 \fillin[{$y = \pm \frac{4}{3}x$}]。
\end{question}

\begin{question}
  若 $(2x - 1)^5 = a_0 + a_1 x + a_2 x^2 + \cdots + a_5 x^5$，则 $a_0 + a_1 + a_2 + \cdots + a_5 = $ \fillin[1]。
\end{question}


% ===== 三、解答题 =====
\section{解答题（本题共 3 小题，共 40 分。解答应写出文字说明、证明过程或演算步骤）}
\examsetup{question/index = 0}

\begin{problem}[points = 12]
  已知函数 $f(x) = 2\sin x \cos x + 2\sqrt{3}\cos^2 x - \sqrt{3}$。

  \begin{enumerate}
    \item 求 $f(x)$ 的最小正周期；
    \item 求 $f(x)$ 在区间 $[0, \frac{\uppi}{2}]$ 上的最大值和最小值。
  \end{enumerate}

  \begin{solution}
    \begin{enumerate}
      \item $f(x) = \sin 2x + \sqrt{3}(2\cos^2 x - 1)$ \score{2}

      $= \sin 2x + \sqrt{3}\cos 2x$ \score{2}

      $= 2\sin(2x + \frac{\uppi}{3})$。\score{2}

      所以最小正周期 $T = \dfrac{2\uppi}{2} = \uppi$。\score{2}

      \item 当 $x \in [0, \frac{\uppi}{2}]$ 时，$2x + \frac{\uppi}{3} \in [\frac{\uppi}{3}, \frac{4\uppi}{3}]$，\score{2}

      当 $2x + \frac{\uppi}{3} = \frac{\uppi}{2}$，即 $x = \frac{\uppi}{12}$ 时，$f(x)$ 取最大值 $2$；\score{1}

      当 $2x + \frac{\uppi}{3} = \frac{4\uppi}{3}$，即 $x = \frac{\uppi}{2}$ 时，$f(x)$ 取最小值 $-\sqrt{3}$。\score{1}
    \end{enumerate}
  \end{solution}
\end{problem}

\begin{problem}[points = 14]
  在 $\triangle ABC$ 中，角 $A$，$B$，$C$ 的对边分别为 $a$，$b$，$c$，已知 $a = 3$，$b = 2\sqrt{3}$，$\cos C = \frac{1}{3}$。

  \begin{enumerate}
    \item 求边 $c$ 的长；
    \item 求 $\triangle ABC$ 的面积；
    \item 求 $\sin A$ 的值。
  \end{enumerate}

  \begin{solution}
    \begin{enumerate}
      \item 由余弦定理 $c^2 = a^2 + b^2 - 2ab\cos C$，\score{2}

      $c^2 = 9 + 12 - 2 \times 3 \times 2\sqrt{3} \times \frac{1}{3} = 21 - 4\sqrt{3}$，\score{2}

      所以 $c = \sqrt{21 - 4\sqrt{3}}$。\score{1}

      \item 由 $\cos C = \frac{1}{3}$ 得 $\sin C = \sqrt{1 - \cos^2 C} = \frac{2\sqrt{2}}{3}$，\score{2}

      $S_{\triangle ABC} = \frac{1}{2}ab\sin C = \frac{1}{2} \times 3 \times 2\sqrt{3} \times \frac{2\sqrt{2}}{3} = 2\sqrt{6}$。\score{2}

      \item 由正弦定理 $\dfrac{a}{\sin A} = \dfrac{c}{\sin C}$，\score{2}

      $\sin A = \dfrac{a\sin C}{c} = \dfrac{3 \times \frac{2\sqrt{2}}{3}}{\sqrt{21 - 4\sqrt{3}}} = \dfrac{2\sqrt{2}}{\sqrt{21 - 4\sqrt{3}}}$。\score{3}
    \end{enumerate}
  \end{solution}
\end{problem}

\begin{problem}[points = 14]
  已知函数 $f(x) = x^3 - 3ax + 2$（$a > 0$）。

  \begin{enumerate}
    \item 讨论 $f(x)$ 的单调性；
    \item 若 $f(x)$ 有三个零点，求 $a$ 的取值范围；
    \item 当 $a = 1$ 时，求 $f(x)$ 在区间 $[-2, 2]$ 上的最大值。
  \end{enumerate}

  \begin{solution}
    \begin{enumerate}
      \item $f'(x) = 3x^2 - 3a = 3(x^2 - a)$。\score{2}

      当 $x < -\sqrt{a}$ 或 $x > \sqrt{a}$ 时，$f'(x) > 0$；

      当 $-\sqrt{a} < x < \sqrt{a}$ 时，$f'(x) < 0$。\score{2}

      所以 $f(x)$ 的递增区间为 $(-\infty, -\sqrt{a})$ 和 $(\sqrt{a}, +\infty)$，

      递减区间为 $(-\sqrt{a}, \sqrt{a})$。\score{2}

      \item $f(x)$ 的极大值为 $f(-\sqrt{a}) = 2 + 2a\sqrt{a}$，

      极小值为 $f(\sqrt{a}) = 2 - 2a\sqrt{a}$。\score{2}

      $f(x)$ 有三个零点等价于 $f(-\sqrt{a}) > 0$ 且 $f(\sqrt{a}) < 0$，\score{2}

      即 $2 - 2a\sqrt{a} < 0$，解得 $a > 1$。\score{1}

      \item 当 $a = 1$ 时，$f(x) = x^3 - 3x + 2$，$f'(x) = 3(x^2 - 1)$。\score{1}

      $f(-2) = -4$，$f(-1) = 4$，$f(1) = 0$，$f(2) = 4$。\score{1}

      所以最大值为 $4$。\score{1}
    \end{enumerate}
  \end{solution}
\end{problem}


% ===== 草稿纸 =====
\newpage
\section*{草稿纸}

\draftpaper

\end{document}
